\section{Can Either Premise Be Checked From Published Papers?}
\label{sec:bg-survey}

The two premises can be checked from a paper only if it reports the per-system scores behind its pooled comparisons, releases enough to attach an uncertainty to them, and states how its scores were produced.
We read \SurveyN microservice RCA method and benchmark papers from 2021 to 2026, a convenience sample built from the benchmark papers, the method papers they compare, and the papers citing them for which we obtained the full text (\cref{tab:survey}), and checked these three things.
A paper counts as reporting per-system results if it gives a score for each evaluated system or dataset.
A release counts as containing per-case outputs if a file in the linked repository or archive gives a method's ranking or prediction for an identifiable case, and as complete if such files exist for every method on every dataset of the paper's main comparison.
Scope is reported: \SurveyMulti of the \SurveyN papers evaluate more than one system, and all \SurveyMulti give a score per system, so a reader can see where a pooled lead does not hold, but not with what uncertainty.
Checking the repositories and archives linked from each paper on 2 and 3 October 2026, we found \SurveyRelAny papers whose release contains per-case outputs of at least one method.
For \SurveyRelDemo of them (\baro~\citep{pham2024baro}, \rcaeval~\citep{pham2024rcaeval}, and a study that shares \rcaeval's repository~\citep{pham2024howfar}) the outputs are those a tutorial notebook saved for a single case.
A further \SurveyRelOwn release the outputs of the authors' own method on their test cases: one archived run of \openrca's agent~\citep{xu2025openrca}, the rankings of CausalRCA~\citep{causalrca}, and the predictions of CHASE~\citep{zhao2024chase}.
One, \circa~\citep{li2022circa}, releases the rankings of every method of its simulation study, but not of its production dataset.
For \SurveyRelUndet papers a linked release had expired or could not be reached.
None of the \SurveyRelFullCheckable papers whose releases we could check releases per-case outputs for every method on every dataset of its main comparison, so for those comparisons a reader cannot compute a paired difference, its uncertainty, or a coverage figure.
\begin{table}[t]
\caption{The \SurveyN surveyed papers. Per-system: a score per evaluated system is reported (single system: one system evaluated). Per-case outputs released: what the release linked from the paper contains, checked on 2 and 3 October 2026. All methods: whether per-case outputs exist for every method on every dataset of the main comparison (partial: for every method of one experiment).}
\label{tab:survey}
\footnotesize
\setlength{\tabcolsep}{3pt}
\begin{tabular}{@{}lllllc@{}}
\toprule
Paper & Year & Venue & Per-system & Per-case outputs released & All methods \\
\midrule
Luan Pham et al. & 2024 & FSE 2024 & yes & one demonstration case & no \\
Mingjie Li et al. & 2022 & KDD 2022 & yes & all methods, simulation study only & partial \\
Azam Ikram et al. & 2022 & NeurIPS 2022 & yes & -- & no \\
Zeyan Li et al. & 2021 & IWQoS 2021 & yes & release unreachable & -- \\
Luan Pham et al. & 2025 & WWW 2025 Companion & single system & one demonstration case & no \\
Aoyang Fang et al. & 2026 & FSE 2026 & yes & -- & no \\
Junjielong Xu et al. & 2025 & ICLR 2025 & yes & own method, all test cases & no \\
Mila Hardt et al. & 2024 & CLeaR 2024 & single system & -- & no \\
Cheryl Lee et al. & 2023 & ICSE 2023 & yes & -- & no \\
Wei Zhang et al. & 2024 & EMNLP 2024 Findings & yes & -- & no \\
Lecheng Zheng et al. & 2024 & WWW 2024 & yes & -- & no \\
Lecheng Zheng et al. & 2025 & IEEE BigData 2025 & yes & -- & no \\
Shuaiyu Xie et al. & 2025 & TOSEM & yes & -- & no \\
Ziming Zhao et al. & 2024 & arXiv 2024 & yes & own method, all test cases & no \\
Luan Pham et al. & 2024 & ASE 2024 & yes & one demonstration case & no \\
Lecheng Zheng et al. & 2024 & OpenReview 2024 & yes & -- & no \\
Zhenhe Yao et al. & 2024 & ISSRE 2024 & single system & -- & no \\
Chenxi Zhang et al. & 2024 & ICSE 2024 & single system & -- & no \\
Ruomeng Ding et al. & 2023 & ESEC/FSE 2023 & single system & -- & no \\
Cheng-Ming Lin et al. & 2024 & AAAI 2024 & yes & -- & no \\
Songhan Zhang et al. & 2025 & arXiv 2025 & yes & -- & no \\
Ruyue Xin et al. & 2023 & JSS 2023 & single system & own method, all test cases & no \\
Shenglin Zhang et al. & 2025 & ISSRE 2025 & yes & release unreachable & -- \\
\bottomrule
\end{tabular}

\end{table}

Whether a number was computed on the raw or the reduced table is a question that even the benchmark's own paper does not answer.
\rcaeval's paper describes the file format of its release but not which of the two tables its reported numbers use, and \cref{sec:rq2-input} shows that the two rank the methods differently.
Superiority claims are nonetheless stated over the benchmark as a whole, for example ``The top-1 accuracy of TraceRCA outperforms the state-of-the-art unsupervised approaches by 44.8\%''~\citep{li2021tracerca} and ``Extensive experiments conducted on the synthetic and realworld microservice-based datasets demonstrate that RUN noticeably outperforms the state-of-the-art root cause analysis methods''~\citep{lin2024run}.
